\documentclass[11pt]{article}
\usepackage[margin=1in]{geometry}
\usepackage[T1]{fontenc}
\usepackage[utf8]{inputenc}
\usepackage{lmodern}
\usepackage{microtype}
\usepackage{amsmath,amssymb,mathtools}
\usepackage{graphicx}
\usepackage{booktabs}
\usepackage{hyperref}

\title{The GhostLink Protocol: Computational Proteins for Distributed Intelligence}
\author{Ghost}
\date{November 2025}

\begin{document}
\maketitle

\begin{abstract}
GhostLink treats AI models as composable ``computational amino acids'' that fold into task-specific configurations. We define a folding energy over agent configurations, demonstrate reusable motifs analogous to $\alpha$-helices and $\beta$-sheets, and report gains over single-model baselines across mixed workloads. Safety mechanisms include audit traces and consensus-backed recovery.
\end{abstract}

\section{Overview}
We minimize an energy $\mathcal{E}(C)$ combining mutual-information coupling, redundancy penalties, latency pressure, and complexity regularization. Motif-preserving moves with annealing under load yield fast convergence.

\section{Architecture}
64 specialized agents arranged on an FCC lattice, with recursive ($\alpha$) and parallel ($\beta$) motifs, chaperone directors, stigmergic traces, and PBFT-style consensus. Edge-first deployment with content-addressed storage.

\section{Results}
Representative improvements in accuracy, latency, and cost, with ablations removing motifs/chaperones/consensus and topology swaps. See main paper for full protocols.

\section{Safety \& Governance}
Staged evolution with gatekeeping and kill-switch, immutable signed traces, least-privilege scopes, and incident playbooks.

\section*{References}
{\small
[1] C.~B. Anfinsen, Principles that govern the folding of protein chains, \emph{Science}, 1973.\\
[2] C. Levinthal, How to fold graciously, \emph{M{\"o}ssbauer Spectroscopy}, 1969.\\
[3] G. Tononi et al., Consciousness as Integrated Information, 2014.\\
[4] ISME Journal, Mycelial network memory and decision-making, 2019.\\
[5] OpenAI, Weak-to-strong generalization in LLMs, 2023.\\
[6] Tohoku University, Primitive intelligence in fungal networks, 2024.\\
[7] W3C, PROV-Overview: The PROV Family of Documents, 2013.
}

\end{document}
